\documentclass[10pt,a4paper]{amsart}
\usepackage[margin=19mm]{geometry}
\usepackage{amsmath,amssymb,mathtools}
\usepackage[scr=rsfs,cal=euler]{mathalfa}
\usepackage{xfrac}
\usepackage[unicode,hidelinks,bookmarks=false]{hyperref}
\newcommand{\ie}{i.e.\ }
\newcommand{\cf}{cf.\ }
\newcommand{\resp}{respectively\ }
\newcommand{\Subsection}{\subsection}
\providecommand{\nobreakdash}{\nobreak\hbox{-}}
\setlength{\emergencystretch}{2em}
\multlinegap0pt
\usepackage{siunitx}
\usepackage{siunitx}
\usepackage{siunitx}
\usepackage{siunitx}
\usepackage{xcolor}
\usepackage{amssymb}
\usepackage{siunitx}
\usepackage{algorithm}\usepackage{algpseudocode}
\usepackage{tcolorbox}
\newtheorem{lemma}{Lemma}
\newcommand{\bv}{\mathbf{v}}
\title{A Discontinuous Galerkin discretization for the sea ice dynamics}
\author{Emma Lagracie, Thomas Richter}
\date{Source: arXiv:2608.26810v1 (CC BY 4.0)}
\begin{document}
\maketitle

\noindent\emph{Source and licence.} A Discontinuous Galerkin discretization for the sea ice dynamics. Version arXiv:2608.26810v1 (CC BY 4.0), distributed by arXiv under the Creative Commons Attribution 4.0 International licence, http://creativecommons.org/licenses/by/4.0/, as recorded in the arXiv licence metadata for this exact version and shown on the abstract page https://arxiv.org/abs/2608.26810v1. Full licence notice: \texttt{licenses/arxiv-2608.26810-CC-BY-4.0.txt}.\par\smallskip

\noindent\textit{Editorial scope.} Counted unit: the article's lemma Bregman (source0.tex lines 1180--1190). The article's complete original proof of it is delivered with it (source0.tex lines 1191--1256, one \texttt{proof} environment of 66 lines). No mathematics is written, repaired or re-derived: the statement and the proof are the article's own lines, in the article's own notation, and no retained line is substituted or re-flowed.  No further source-local context is needed: every cross-reference of every retained line resolves inside the two delivered spans.  Nothing was omitted from any retained span, so the package carries no omission record.  No retained line contains a citation, so the wrapper carries no bibliography and no bibliography entry.  No retained line contains a figure environment, an image-inclusion command or drawing code, so no image is delivered and the package declares no asset dependency.  Editorial formatting only, in addition: an AMS \texttt{amsart} preamble that restates the article's own declarations and macros used by the retained lines, namely the environments \texttt{lemma} on separate counters, together with the standard packages those lines need.  The retained environments print in this document as Lemma 1 in order of appearance, whereas the article prints the same items with the numbering of its own full document.  The complete native source of the article as distributed in the arXiv e-print for this exact version is bundled unchanged as \texttt{sources/208663/source0.tex}.
\begin{lemma}[Bound for the Bregman divergence]
\label{lemma: bound Bregman}
Let $\psi(\bv)=\frac13 \|\bv\|_2^3.$
Then, for all \(\bv,\bv^*\in\mathbb R^2\) we have the following pointwise inequality
\[
    D_\psi(\bv,\bv^*)
    \geq
    \frac1{12}
    \left (\|\bv\|_2+\|\bv^*\|_2\right )\|\bv -\bv^*\|_2^2.
\]
\end{lemma}

\begin{proof}
Define the pointwise Bregman divergence associated with \(\psi\) by $ D_\psi(\bv,\bv^*)
    =
    \psi(\bv)-\psi(\bv^*)-\nabla\psi(\bv^*)\cdot(\bv-\bv^*).$
By Taylor's formula with integral remainder, and taking $h=\bv-\bv^*$, we have
\[
    D_\psi(\bv,\bv^*)
    =
    \int_0^1
    (1-t)\,
    h^T\nabla^2\psi(\bv^*+th)h\,dt.
\]
Since $\nabla^2\psi(\bv)=
    \|\bv\|_2 I+\frac{\bv\otimes \bv}{\|\bv\|_2} \succeq
    \|\bv\|_2 I$, we get
$$
    D_\psi(\bv,\bv^*)\geq \|h\|_2^2 \int_0^1
    (1-t)\|\bv^*+t h\|_2\,dt.
$$
Moreover, $\bv^*+t h=(1-t)\bv^*+t \bv$, then,
\[
\|(1-t)\bv^*+t\bv\|_2 \geq \bigl| (1-t) \|\bv^*\|_2-t\|\bv\|_2 \bigr|.
\]
Therefore,
\[
    D_\psi(\bv,\bv^*)
    \geq
    \|h\|_2^2 \int_0^1 (1-t)
    \bigl| (1-t)\|\bv^*\|_2-t\|\bv\|_2 \bigr|\,dt.
\]
We set $a=\|\bv^*\|_2\ge 0$ and $b=\| \bv\|_2 \ge0$.
If \(a=b=0\), then \(\bv=\bv^*=0\), and the result is immediate. Otherwise,
\[
    \bigl| (1-t)a-tb \bigr|
    =
    (a+b)\left| \frac{a}{a+b}-t\right|,
\]
hence we have
\[
    D_\psi(\bv,\bv^*)
    \geq
    (a+b)\|h\|_2^2
    \int_0^1
    (1-t)\left |\frac{a}{a+b}-t\right |\,dt.
\]
We can compute the integral to get
\[
    \int_0^1
    (1-t)\left |\frac{a}{a+b}-t\right |\,dt
    =
    -\frac{1}{3}\left (\frac{a}{a+b}\right )^3
    +\left (\frac{a}{a+b}\right )^2
    -\frac{1}{2} \left (\frac{a}{a+b}\right )
    +\frac16.
\]
The function $f: x \mapsto f(x) = -\frac{1}{3}x^3+x^2-\frac{1}{2} x+\frac16$ is bounded from below on \( [0,1]\). On this interval, its local extrema are localized in $0,\, 1$ and $1- \frac{\sqrt{2}}{2}$ with
\[f(x) = \int_0^1(1-t)|x-t|\,dt \geq \frac1{12},\quad \forall x\in [0, 1].
\]
Finally, we have the lower bound
\[
    D_\psi(\bv,\bv^*)
    \geq
    \frac1{12}
    \left (\|\bv\|_2+\|\bv^*\|_2\right )\|\bv -\bv^*\|_2^2.
\]
\end{proof}

\end{document}
